\section{Traversal AST dan Visitor di C}
\label{sec:ch16-traversal}

\subsection{Traversal Rekursif}

Traversal AST adalah proses mengunjungi simpul-simpulnya. Dua urutan
umum:
\begin{itemize}[leftmargin=*]
  \item \textbf{pre-order}: kunjungi simpul, lalu anak-anak;
  \item \textbf{post-order}: kunjungi anak-anak, lalu simpul.
\end{itemize}

Pre-order cocok untuk analisis yang membutuhkan konteks dari orang
tua, misalnya validasi loop. Post-order cocok untuk sintesis,
misalnya perhitungan tipe.

\begin{figure}[H]
\centering
\begin{tikzpicture}[node distance=6mm and 10mm]
  \node[phase, text width=2.6cm] (root) {\code{+}};
  \node[phase, text width=2.0cm, below left=of root] (a) {\code{a}};
  \node[phase, text width=2.0cm, below right=of root] (mul) {\code{*}};
  \node[phase, text width=2.0cm, below left=of mul] (b) {\code{b}};
  \node[phase, text width=2.0cm, below right=of mul] (c) {\code{c}};
  \draw[arrow] (root) -- (a);
  \draw[arrow] (root) -- (mul);
  \draw[arrow] (mul) -- (b);
  \draw[arrow] (mul) -- (c);
\end{tikzpicture}
\caption{Urutan pre-order: \(+\), \(a\), \(*\), \(b\), \(c\)}
\label{fig:ch16-preorder}
\end{figure}

\subsection{Pola Visitor}

Pola visitor memisahkan operasi dari struktur data. Pada C, visitor
dapat diimplementasikan dengan:
\begin{itemize}[leftmargin=*]
  \item fungsi yang menerima simpul dan konteks;
  \item tabel fungsi per jenis simpul;
  \item \code{switch} pada \code{kind}.
\end{itemize}

Proyek Subset-C memakai \code{switch} karena sederhana dan mudah
dibaca. Setiap operasi baru ditambahkan sebagai fungsi traversal
terpisah, misalnya \code{typecheck} dan \code{validate}.

\subsection{Contoh Traversal}

Fungsi \code{visit\_children} mengunjungi semua anak:
\begin{lstlisting}[language=C]
static void visit_children(Context *context, AstNode *node) {
    AstNode *child;
    for (child = node->children; child; child = child->next) {
        visit(context, child);
    }
}
\end{lstlisting}
Fungsi \code{visit} menangani jenis simpul dan memanggil
\code{visit\_children} bila perlu. Pola ini menghindari duplikasi
kode traversal.

\subsection{Konteks Traversal}

Traversal sering membutuhkan konteks, misalnya kedalaman loop atau
tipe kembalian fungsi. Konteks diteruskan sebagai parameter dan
diperbarui saat masuk/keluar simpul. Pada validator, konteks berisi
\code{loop\_depth}, \code{return\_type}, dan \code{has\_return}.
